\scanpage{45}
Il y a correspondance biunivoque entre \uline{ces} sous-tores et les sections de $\mathrm{Kill}/(\mathbb Z/2\mathbb Z)$ [NB ``\uline{le}'' $W$ opère sur $\mathrm{Kill}$, et dans $W$ il y a ``\uline{le}'' élément $\underline w$ d'ordre $2$ qui transforme \uline{la} \deleted{B} chambre de Weyl\inserted{\unsure{standart|standard}} en son opposée, c'est par ce $\underline w$ qu'on divise \ldots] : couples de un tore maximal et d'un couple de Borel opposés le contenant (mais ces Borels n'étant pas néc. définis séparément sur $S$).

\margin{\unsure{si|G} simple ?}
(ii) Soit $i:G_{\mathrm m}\unsure{-|\to}G$ un hom principal. Il induit un hom
\[
\unsure{\kappa|K}_{\deleted{\unsure{h}}}:\mu_h\longrightarrow G.\qquad h\ \text{le nb de Coxeter}
\]
Un hom $\kappa_{\deleted{\unsure{h}}}$ est dit \deleted{\unsure{de}}\inserted{\uline{hom}} \uline{de Kostant} si loc fpqc (ou ét.f., cela vient au $=$, comme on voit aussitôt) il est obtenu à partir d'un hom principal.

\deleted{Le} Alors $\kappa_{\deleted{\unsure{h}}}$ est régulier, i.e. le centralisateur connexe de $\kappa_h$ est un tore maximal. Faisant opérer $\mu_h$ sur $\mathfrak g\mathrel{\underset{\unsure{*}}{=}}\underline{\mathrm{Lie}}(G)$, la graduation de type $\mathbb Z/h\mathbb Z$ correspondante de $\mathfrak g$ donne comme composante de degré $0$ une \deleted{\unsure{tore maximal} $T=$ \unsure{Zer}}\inserted{sous-Algèbre de} Cartan $\mathfrak t=\deleted{\unsure{\mathfrak g^0}}\mathrm{Lie}(T)$, $T=\underline{\mathrm{Cent}}^{0}(i_h)$, et la composante de degré $j$, $1\leq j\leq h-1$, \deleted{est égal au}\inserted{a comme} rang le nombre de racines positives dont l'ordre est $\equiv j\ (h)$, i.e. \unsure{don|dont} l'ordre est $j$ ou $j-h$ ; en fait, si $\kappa$ est associé à un couple de Killing, \deleted{\unsure{sur}}, \deleted{déployé} avec syst. de racines associé à $T$ \uline{déployé}, alors
% archive sheet number: 44

